% Anschreiben: {{EMPLOYER}}, {{ROLE}}
% Erzeugt am {{DATE}} von apply/new.py
%
% Nur der FIT-Block muss geschrieben werden. Der Rest stimmt und ist wiederverwendbar.
% Kompilieren mit pdflatex lokal oder auf Overleaf.

\documentclass[a4paper,11pt,version=last]{scrlttr2}
\usepackage[utf8]{inputenc}
\usepackage[T1]{fontenc}
\usepackage[ngerman]{babel}
\usepackage{lmodern}
\usepackage[colorlinks=true, urlcolor=blue]{hyperref}

\setkomavar{fromname}{Baha Khammari}
\setkomavar{fromaddress}{D\"usseldorf, Deutschland}
\setkomavar{fromemail}{b.e.khammari@gmail.com}
\setkomavar{fromphone}{+49\,1520\,9016956}
\setkomavar{subject}{Bewerbung als {{ROLE}}}
\KOMAoptions{fromemail=true, fromphone=true, backaddress=false}

\begin{document}

\begin{letter}{%
    {{CONTACT}} \\
    {{EMPLOYER}} \\
    {{LOCATION}}
}

\opening{{{OPENING}},}

% ---------------------------------------------------------------------------
% FIT. Der einzige wirklich neue Teil. Zwei bis drei S\"atze, die auf keinen
% anderen Arbeitgeber passen w\"urden. Etwas Konkretes aus der Ausschreibung
% aufgreifen. L\"asst sich das nicht schreiben, passt die Stelle vermutlich nicht.
%
% Muster:
%   Ihr Team arbeitet an [konkret aus der Ausschreibung].
%   In meiner Masterarbeit habe ich [Methode oder Thema], und
%   [eine konkrete geforderte F\"ahigkeit].
%   Genau das m\"ochte ich bei [das, was sie nennen] einbringen.
% ---------------------------------------------------------------------------

{{FIT}}

% --- Belege. Wiederverwendbar. Auf eine Seite k\"urzen. ---
In meiner Masterarbeit (Note 1,0) habe ich die kausalen Lohneffekte des brasilianischen Stipendienprogramms ProUni gesch\"atzt. Daf\"ur habe ich in \textbf{R} rund f\"unf Millionen Arbeitgeber-Arbeitnehmer-Datens\"atze aus 558 Regionen aufbereitet und mit einem heterogenit\"atsrobusten Differenz-von-Differenzen-Sch\"atzer ausgewertet. Das Ergebnis ist ein nicht-monotoner Effekt, dessen Vorzeichen g\"angige Verfahren umkehren. Solche Befunde h\"alt nur sauberes Arbeiten mit gro\ss{}en Verwaltungsdaten aus.

Die n\"otige Disziplin bringe ich aus meiner Zeit in Assurance Solutions bei \textbf{PwC Deutschland} mit. Dort war Excel das t\"agliche Werkzeug, um Kennzahlen zu pr\"ufen und Ergebnisse f\"ur das Management aufzubereiten. Ich wei\ss{} also, wie wichtig Pr\"azision ist, wenn andere auf die Zahlen bauen, und wie man aus einem komplexen Datensatz eine klare Aussage macht.

Mein Masterstudium schlie\ss{}e ich im September 2026 ab, ab Oktober 2026 bin ich voll verf\"ugbar.

\"Uber ein Gespr\"ach w\"urde ich mich sehr freuen.

\closing{Mit freundlichen Gr\"u\ss{}en}

\end{letter}

\end{document}
